\newpage
\section{M.Carlsen - W.Anand, World Chess Championship}
\begin{multicols}{2}
    \newchessgame
    \chessgameinfo{World Chess Championship Game 2}{M.Carlsen}{W.Anand}{}{2014}{1-0}
    \mainline[level=1]{
        1. e4 e5 2. Nf3 Nc6 3. Bb5 Nf6 4. d3 Bc5 5. O-O d6 6. Re1 O-O
    }

    \mythoughts{
        The position resembles the Italian or the Ruy Lopez, where White often reroutes
        the b1-knight to f5 via g3 or e3 and builds a center with c3 and d4.
        Compared with those standard structures, Black's dark-squared bishop is on c5
        rather than e7. I would still choose \variation{7. c3}. If
        \variation{7... a6 8. Ba4 b5 9. Bc2}, then once I establish the center, Black's c7-pawn
        can become a target.
    }

    \myreflection{
        My move is playable, but the claim that ``Black's c7-pawn can become a target'' is wrong:
        the bishop can simply protect it from b6.
    }

    \mainline[level=1]{
        7. Bxc6 bxc6
    }

    \mythoughts{
        Candidate moves are \symknight bd2, c3, and h3. \variation{8. h3} is less relevant,
        because Black is unlikely to play \symbishop g4 and exchange that knight, as it
        would help White open the g-file. I would choose \variation{8. Nbd2}.
    }

    \myreflection{
        My move is not good.
        True, Black is unlikely to play \symbishop g4, but he can play \symknight g4,
        so h3 does have its merits.
        Also, \variation[invar]{8. c3 Ng4 9. d4} is playable.
    }

    \mainline[level=1]{
        8. h3 Re8
    }

    \mythoughts{
        The usual plan for building the center is \symknight bd2, c3, \symqueen c2, and then d4.
        I still prefer \variation{9. Nbd2}.
    }

    \mainline[level=1]{
        9. Nbd2 Nd7
    }

    \chessboard

    \mythoughts{
        Black wants to play \symknight f8-e6 to stop me from establishing a pawn center.
        If I fail to set up that center, I do not see how to improve my position.
        So I need to hurry with \variation{10. c3}.
    }

    \myreflection{
        My move is playable. Once I can match the position to something in my knowledge base,
        I can play a series of moves without calculating each one.

        Magnus Carlsen thinks only when he cannot find a similar position;
        otherwise he simply reuses an existing solution.
    }

    \mainline[level=1]{
        10. Nc4 Bb6
    }

    \mythoughts{
        I want to place my dark-squared bishop on b2, as that seems to be its best square.
        I also want to play g3, \symknight h4, f4, \symknight f5 etc. 
        Black must watch out for sacrifices on e5 and pressure against g7.
        
        After \variation{11. a4 a5}, I don't get anything since I have no access to b5.
        
        I want to play \variation{11. b3}.
    }

    \myreflection{
        My claim that ``after \variation{11. a4 a5}, I don't get anything since I have no access to b5''
        is wrong. After the exchange on b6, Black can no longer recapture with the a-pawn,
        which opens the a-file. Such an exchange therefore favors White.

        My move is playable.
        However, the plan of placing the bishop on b2 is not realistic: White
        must reckon with Black's central strike d5:
        \variation[invar]{
            11. b3 d5 12. exd5 cxd5 13. Ncxe5 Nxe5 14. Nxe5 Rxe5 15. Rxe5 Bxf2+ 16. Kh2 Bd4 17. Bf4 Bxe5 18. Bxe5 
        }
        The line is almost forced, and I had not calculated it at all.
    }

    \mainline[level=1]{
        11. a4 a5
    }

    \chessboard

    \mythoughts{
        Playing for an attack on the kingside is a typical
        motif for Ruy Lopez structures. One idea would be 
        g3, h4, \symking g2, \symrook h1. However, I must account for
        Black can strike in the center with d5, e4, \symknight e5 etc. 
        The attack is stopped.

        The other route is \symknight g5, h4, \symrook e3. To make this work,
        I first need to remove the bishop on b6. Therefore \variation{12. Nxb6}.
    }

    \mainline[level=1]{
        12. Nxb6 cxb6
    }

    \chessboard
    
    \mythoughts{
        I do not like attacking with f4 yet because my rooks are not ready.
        Black can defend f7 without difficulty.

        I would like to attack with \symknight g5. In \variation{13. Ng5 h6 14. Nxf7 Kxf7 15. Qh5+ g6}, I
        don't see how to proceed. \variation{13. Ng5 h6 14. Nxf7 Kxf7 15. Bxh6 gxh6 16. Qh5+ Ke7 17. Re3 Nf6}
        I still do not see a clear continuation. My pieces are not fully developed, and sacrificing
        two pieces is too early.

        \variation{13. h4} looks better. Black would prefer not to play h6 first because it weakens the kingside.
        \variation{13. h4 Nf6 14. Bg5} the pin is unpleasant for Black. 
        \variation{13. h4 d5} is unpleasant for White. White must wait for Black to play dxe4. Black
        can then move his knight to e6 targeting d4. c3 would not help because it weakens the d3 pawn 
        permanently.

        Clearly I need another idea. I wanted to play d4 anyway, so what about doing it now?
        \variation{13. d4 exd4 14. Nxd4} hitting c6 and f5. It looks good for me.
        I also intend to play \variation{14. dxe5}, \variation{15. Be3} targeting the b6 pawn.
        
    }

    \myreflection{
        The variations above are based on a single idea and should be easy to defend.
        In practice, I should not spend much time calculating such lines, because it is
        unlikely that I will get much from them.

        I was mixing ideas from the game Carlsen vs. Wang Hao, Tata Steel 2011.

        \chessboard{
            setfen=2rqkb1r/1p1n1pp1/p1n1p1p1/2ppP3/3P4/2P1BN2/PP2BPPP/R2Q1RK1 w k - 0 13
        }

        There, the center is completely closed, so the pawn thrusts g3 and h4 are effective
        because Black cannot cover the weakness with his rooks.

        In the actual game, even if I open lines successfully, Black can simply place
        his rooks on the 7th rank and my attack is over.

        A more subtle approach is to focus on the opponent's weaknesses:
        tie him down to defending one weakness, then strike somewhere else where the defense
        is thin. Alekhine won many games with this idea.

        As for the concrete move, after \variation[invar]{13. h4 c5} Black is better, because White
        can never set up the center. The lesson here, clichéd as it is: pay attention
        to a central counterstrike whenever you plan to play on the flank.
    }

    \mainline[level=1]{
        13. d4 Qc7
    }

    \chessboard 

    \mythoughts{
        Black's last move protects the weak pawns on the queenside. 

        My c1 bishop is not active yet, and I am not sure where to place it. Both a3 and e3 are
        reasonable squares for it. 
        My a1 rook is the worst piece and therefore I want to improve it.
        Since I do not yet know where to develop the c1 bishop, I do not want to keep
        the rook on the first rank; it should wait until the bishop is developed.

        Developing the rook via a3 looks good. I can then place it on d3 and target d6. My
        bishop can either go to a3, e3 or even g5 depending on what happens next.

        I do not see any immediate threat from Black. \variation{14... d5} drops a pawn after
        \variation{15. dxe5 Nxe5 16. exd5 exd5 17. Qxd5} 

        So I decide to play \variation{14. Ra3}.
    }

    \myreflection{
        Yes, I found the same move as Carlsen!

        By the way, the engine recommends \variation[invar]{
            14. c4
        }, which is the same idea as my next move.
    }

    \mainline[level=1]{
        14. Ra3 Nf8
    }

    \mythoughts{
        I would love to place my knight on f5. If Black's bishop takes it,
        I recapture with the e-pawn, and f6 can become a serious threat.

        That idea does not work because Black can take on d4 with exd4.
        I have a weakness on the first rank and could lose a bishop.

        This means the center must be stable before I take any action 
        on the kingside. c4 and then d5 is natural for me.

        Another idea is \symrook d3. After exchanging the e-pawns,
        the heavy pieces will be all exchanged. Do I have advantage with 
        the remaining knight and bishop endgame? Black can place his 
        bishop on e6, denying the c4 square for my knight. We both have 
        a good bishop. I do not see how I can break through, since my king has no route
        through the light squares; they are all covered by
        Black's bishop.

        I decide to play \variation{15. c4}.
    }

    \myreflection{
        \variation[invar]{
            15. Rd3
        } is playable, but White should then close the center with d5.

        After \variation[invar]{
            15. Rd3 Ng6 16. dxe5?! dxe5
        }, Black need not exchange pieces, for the same reason: White's
        heavy pieces have no targets. It is therefore unsound to speculate about an endgame
        that would not actually arise here.
    }

    \mainline[level=1]{
        15. dxe5 dxe5
    }

    \mythoughts{
        I thought I had to place my rook on d3, otherwise Black would occupy
        the d-file. After deeper consideration, I no longer think it is mandatory:
        my first and second ranks are covered. Black's rooks have no squares to invade.

        If I still want to attack on the kingside, I need to move my f3-knight
        so that the a3-rook can swing over.

        \variation{16. Nh4 Rd8 17. Qh5} is natural. All my pieces can join the 
        attack soon.

        There is a clear lesson here: occupying an open file does not always bring
        an advantage when the rooks have no targets.
    }

    \mainline[level=1]{
        16. Nh4 Rd8 17. Qh5 f6
    }

    \mythoughts{
        I realize \variation{18. Nf5? g6} I will lose material.
        \variation{18. Rg3} is natural. Next I can bring my knight to f5. Black still needs to develop the
        bishop. I feel I have some advantage now.
    }

    \myreflection{
        The statement above is wrong:
        \variation[invar]{18. Nf5 g6? 19. Rg3 Kh8 20. Qh4 Bxf5 21. Qxf6+ Qg7 22. Qxg7+ Kxg7 23. exf5 }.

        My intended move is still good, however, since I can follow up with \variation{19. Nf5}.
    }
    \mainline[level=1]{
        18. Nf5! Be6?!
    }

    Incidentally, this is a poor move: the e6 square belongs to the knight, not the bishop.

    \mythoughts{
        I don't need to worry about \variation{19... Bxf5}. In the endgame 
        this bishop is the most important defender. When it gets traded,
        I will have a winning endgame. 

        I still want to play \variation{19. Rg3}.

    }

    \mainline[level=1]{
        19. Rg3 Ng6
    }

    \chessboard

    \mythoughts{
        I still need more pieces for the attack. Concrete
        calculation is not yet necessary. This saves me time and energy.

        Bringing my e1-rook to the g-file does not help much. It costs tempi and
        allows Black to counterattack with \symrook d1+.

        Even without that counterattack, I can have at most four attackers
        against the g7 pawn (two rooks on g-file, one f5 knight and queen or bishop on h6).
        Black already has four defenders. So this brute-force attack does not work
        out.

        The same argument also applies to my c1 bishop, since it can only attack g7 as
        its target. 

        Through this process of elimination, I conclude I must advance the h-pawn. Once it reaches h6,
        Black may be forced to play g6, which allows further piece sacrifice on g6.

        I must move my queen to make room for the h-pawn. Placing the queen on g4 lets her
        protect the h-pawn on h4; otherwise Black can take it by first exchanging the knight
        on f5 and then take on h4.
        
        I decide to play \variation{20. Qg4}.
    }

    \myreflection{
        I like that I am making logical decisions, even if the conclusion is not entirely correct.
        After \variation[invar]{
            20. Qg4 Kh8 21. h4 Bxf5
        } I should recapture with the queen rather than the e-pawn, because otherwise Black
        can play \variation{22... Rd4}.

        After \variation{
            20. Qg4 Kh8 21. h4 Bxf5 22. Qxf5 Nxh4 23. Qh5 g5 24. Bxg5 fxg5 25. Rxg5 Nxg2 26. Rxg2 Rg8 
        }
        the game is equal.
        So here it does not matter whether I place the queen on g4 or f3, as both lead
        to the same variation.
    }

    \mainline[level=1]{
        20. h4 Bxf5
    }

    \mythoughts{
        This is indeed a strange move. Black has already lost a tempo
        with the bishop, and exchanging it for my knight is also a bad idea,
        since I have established that the endgame favors me.

        I decide to play \variation{21. exf5}. It drives the knight away, and a pawn
        on f5 will also help me after I play h6 and Black responds with g6.
    }

    \mainline[level=1]{
        21. exf5 Nf4
    }

    \mythoughts{
        I decide to play \variation{22. Qg4}. I want my queen on the kingside for the attack.
        I don't see any other squares for her.
    }

    \myreflection{
        The queen on h5 is actually well placed: as we will see, it controls
        both the e8 and the d1 squares. I failed to reason soundly here.
        Even though the queen needed a better square, g4 is not it, since
        I had already established that I do not have enough attackers against g7.

        The text move is quite logical: remove the best defender.
        Besides, my bishop has not moved at all, and the knight on f4 cost Black
        a lot of tempi, so I have a clear lead in development.
    }

    \mainline[level=1]{
        22. Bxf4 exf4
    }

    \mythoughts{
        I have calculated \variation{23. Rc3 Rd5 24. Qf3 Qb7 25. Re6 Rc8 26. Rexc6 Rxc6 27. Qd5+} winning material.
        I don't like \variation{23. Rc3 c5} because my rook is not doing
        a lot on c3. 

        I don't like \variation{23. Rg4} at all because my queen is stuck on h5.

        I do not like \variation{23. Rf3} either. I need more tempi to win the f4-pawn.
        My pieces are cramped. It does not look
        like a good deal.

        I don't like \variation{23. Rg6} because the move is too committal: my queen and rook
        are stuck on the kingside and I don't see how to proceed. 

        I must play \variation{23. Rc3}; it is the least painful option.
    }

    \myreflection{
        \variation[invar]{23. Rc3 Rd5 24. Re8 Rxe8 25. Qe8#} mate! As we will see,
        this is not the only time I miss this simple tactic.

        \variation[invar]{23. Rf3 }
    }

    \mainline[level=1]{
        23. Rc3 c5
    }

    \chessboard

    \mythoughts{
        My queen on h5 doesn't do much and therefore I want to improve it. 
        \variation{24. Qf3} is a natural move. \variation{25. Rc4} will be the next move, 
        winning the f4 pawn.
    }

    \myreflection{
        The f4-pawn is not important at all. The queen on f3 is worse than on h5,
        as we have already established.
    }

    \mainline[level=1]{
        24. Re6 Rab8
    }

    \mythoughts{
        \variation{25. Qe2} threatens \variation{26. Re7}, which is a natural follow-up.
    }

    \mainline[level=1]{
        25.  Rc4 Qd7
    }

    \chessboard

    \mythoughts{
        Black threatens to play \variation{26... Qd1 27. Qxd1 Rxd1+ 28. Kh2 Rd2} Black 
        has an advantage in the endgame because the rook on d1 is more active.

        I do not see how \variation{26. Kh2 Qd1+ 27. Qxd1 Rxd1} helps me. My f2-pawn is too weak.
        I want to play \variation{26. Rce4 Qd1+ 27. Qxd1 Rxd1+ 28. Kh2 Rd2 29. R1e2 Rxe2 30. Rxe2}.
        My king arrives in time to protect the f5-pawn and take the f4-pawn. Black's rook
        cannot stay on d4 because I can play b3 and then c3, kick it away.
    }

    \myreflection{
        ``I do not see how \variation[invar]{26. Kh2 Qd1+ 27. Qxd1 Rxd1} helps me.'' It is
        painful to miss a simple mate in two: \variation{27.Re8 Rxe8 28. Qe8#}.

        My calculation is correct, but my evaluation of the rook endgame is also wrong.
        After \variation{26. Rce4 Qd1+ 27. Qxd1 Rxd1+ 28. Kh2 Rd2 29. R1e2 Rxe2 30. Rxe2 b5}, the endgame is equal.
    }

    \mainline[level=1]{
        26. Kh2 Rf8
    }

    \mythoughts{
        I must keep queens on the board. If Black manages to exchange queens
        on d1, his rook on d1 cuts my king off. I don't think I can win this endgame.

        \variation{27. Rce4} prevents Black from playing \variation{27... Rbe8}. I threaten
        \variation{28. Re7} in the next move.
        \variation{27. Rce4 Qc7 28. Re7} is winning for White since White can play \variation{29. Ra7} and 
        then \variation{30. Ree7}.

        \variation{27. Rxf4 } is worse since the f4-pawn is not the focus of the game.

        I decide to play \variation{27. Rce4}.
    }

    \myreflection{
        I am glad to see that I recognized the f4-pawn is not the focus of the game.
    }

    \mainline[level=1]{
        27. Rce4 Rb7
    }

    \mythoughts{
        I want to play \variation{28.Qe2}, threatening \variation{29. Re7 Qc8 30. Rxb7 Qxb7 31. Re7 Qc8 32. Qg4 Rf7 33. Rxf7 Kxf7 34. Qxf4}.
        \variation{29. Re7 Qc8 30. Qg4} is even better because Black cannot answer with \variation{32... Rf7}.
    }

    \mainline[level=1]{
        28. Qe2 b5 }
    
    \mythoughts{
        I do not see what Black is up to. I want to carry out my planned
        \variation{29. Re7}.
    }

    \myreflection{
        I could not calculate the whole line:
        \variation[invar]{
            29. Re7 Qd6 30. f3! Rxe7 31. Rxe7 bxa4 32. Qe4 Qb8 33. Qxa4 Qd6 34. Re4 h6 35. Qxa5 Kh7 
        }
        Here f3 is an important prophylactic move.
    }
    \mainline[level=1]{29. b3 bxa4
    }

    \chessboard

    \mythoughts{
        I will take the pawn with \variation{30. Rxa4}. Allowing Black to play ...a3 
        does not make sense since I do not have a mating attack. Taking with the b-pawn opens
        the b-file.
    }

    \myreflection{
        My move keeps the b-file closed, but misses the more important e-file:
        \variation[invar]{
            30. Rxa4 f3 31. Qxf3 Rb4 32. Ra6!
        }
        White still has enough advantage, but I am not sure I would find the winning line after
        \variation{31... Rb4}.
    }
    \mainline[level=1]{
        30. bxa4 Rb4
    }

    \mythoughts{
        Suddenly I see \variation{31. Re7}. Black must take the a4 pawn otherwise
        \variation{32. Rxb4 axb4 33. Qc4+ Kh8 34. Qf7}.
        After \variation{31... Qxa4 32. Rxg7 Kxg7 33. Qg4+ Kh6}, I don't see a checkmate.

        \variation{32. Qg4 Rf7 33. Re8+ Rf8 34. R4e7} White is winning.
    }

    \mainline[level=1]{
        31. Re7 Qd6
    }

    \mythoughts{
        The original plan doesn't work since Black's queen covers the f8 square. 
        \variation{
            32. Qf3
        } threatens \variation{33. Rxb4} and then \variation{34. Qb7}. 
    }

    \mainline[level=1]{
        32. Qf3 Rxe4 33. Qxe4 f3+ 34. g3 h5 35. Qb7
    }

    \subsection*{My Reflection}
    I am glad I was able to make moves based on logic. This way, I can arrive at
    good decisions without calculating a lot.

    During the game I also learned that occupying an open file is not always a good idea
    when the rooks have no target. With this in mind, I should no longer assume that
    occupying an open file is automatically the right move.

    Missing the back-rank tactic was painful: because of that oversight, I was forced
    to settle for suboptimal moves.


\end{multicols}